\chapter{Statistical Mechanics and the Origins of Learning Machines}
\label{ch:statmech}\label{ch:background}

Statistical mechanics is the discipline that explains how probabilistic
laws emerge from deterministic dynamics. It is also the field in which the
neural network was invented. This chapter serves both facts, and its two
halves have distinct jobs. The first half builds the small set of physical
tools that the rest of the thesis uses, beginning with the principle of
stationary action --- the variational pattern that recurs as free-energy
minimisation, as the training objective of a diffusion model, and as
maximum likelihood --- and then (the Boltzmann distribution, the
partition function and free energy, and the Ising energy), with the
derivations carried out in full: the Boltzmann distribution is the
stationary law of Langevin dynamics (Chapter~\ref{ch:stochproc}), the free
energy is the variational principle behind the diffusion-model training
objective (Chapter~\ref{ch:diffusion}), and sampling by simulating a
physical process is statistical mechanics turned into an algorithm. The
second half is the part of the literature review that most histories of
machine learning skip: how spin glasses became the Hopfield network, how
the Hopfield network became the Boltzmann machine, and how the statistical
mechanics of these systems grew into the modern theory of learning. This
is the tradition, recognised by the 2024 Nobel Prize in Physics, in which
this thesis works. The physical presentation follows classical references
\citep{goldstein2002classical, landau1980statistical}.

\section{The principle of stationary action}\label{sec:action}

Mechanics can be built from a single variational statement, and it is worth
starting there rather than with Hamilton's equations, because the same
statement --- \emph{a physical trajectory is a stationary point of a
functional} --- returns three more times in this thesis: as the free-energy
principle of \S\ref{sec:free-energy}, as the variational bound that trains a
diffusion model (Chapter~\ref{ch:diffusion}), and as the maximum-likelihood
problem solved by expectation--maximisation in
Chapter~\ref{ch:em-method}. The objects change; the shape of the argument
does not.

Let a system be described by generalised coordinates $q(t)$ and let the
\emph{Lagrangian} $L(q, \dot q, t)$ be the difference between kinetic and
potential energy, $L = T - U$. For a path $q(t)$ joining fixed endpoints
$q(t_1) = q_1$ and $q(t_2) = q_2$, define the \emph{action}
\begin{equation}
  S[q] \;=\; \int_{t_1}^{t_2} L\big(q(t), \dot q(t), t\big)\,\dd t .
  \label{eq:action}
\end{equation}
The action assigns one number to a whole path. Hamilton's principle states
that the path actually taken is a \emph{stationary point} of $S$: for any
smooth variation $\delta q(t)$ vanishing at the endpoints,
$\delta S = 0$.

\begin{remark}[Stationary, not minimal]
It is usually called the principle of \emph{least} action, and that is
slightly wrong. The physical path is stationary, not necessarily minimal ---
it can be a saddle. The same caution applies everywhere the pattern recurs:
expectation--maximisation converges to a stationary point of the likelihood,
not to a global maximum (Chapter~\ref{ch:em-method}), and it is a recurring
source of false confidence to read ``the algorithm converged'' as ``the
algorithm found the best answer''.
\end{remark}

Carrying out the variation gives the Euler--Lagrange equations. Writing
$q(t) \to q(t) + \delta q(t)$ and expanding to first order,
\[
  \delta S = \int_{t_1}^{t_2}
  \left(\frac{\partial L}{\partial q}\,\delta q
      + \frac{\partial L}{\partial \dot q}\,\delta \dot q\right)\dd t ,
\]
and integrating the second term by parts, with the boundary term vanishing
because $\delta q(t_1) = \delta q(t_2) = 0$,
\[
  \delta S = \int_{t_1}^{t_2}
  \left(\frac{\partial L}{\partial q}
      - \frac{\dd}{\dd t}\frac{\partial L}{\partial \dot q}\right)\delta q \,\dd t .
\]
Since this must vanish for every admissible $\delta q$, the bracket must
vanish pointwise:
\begin{equation}
  \frac{\dd}{\dd t}\frac{\partial L}{\partial \dot q_i}
  \;-\; \frac{\partial L}{\partial q_i} \;=\; 0,
  \qquad i = 1, \dots, n.
  \label{eq:euler-lagrange}
\end{equation}
For $L = \sum_i \tfrac12 m_i \dot q_i^2 - U(q)$ this returns
$m_i \ddot q_i = -\partial U/\partial q_i$, which is Newton's second law. The
variational principle is therefore not a reformulation for its own sake: it
derives the dynamics from a scalar functional, and that is the move which
generalises.

\paragraph{From Lagrange to Hamilton.} The passage to the Hamiltonian
picture is a Legendre transform. Define the conjugate momentum
$p_i = \partial L/\partial \dot q_i$ and set
\begin{equation}
  H(q,p) \;=\; \sum_i p_i \dot q_i \;-\; L(q,\dot q),
  \label{eq:legendre}
\end{equation}
with $\dot q$ eliminated in favour of $p$. For $L = T - U$ with $T$ quadratic
in the velocities this gives $H = T + U$, the total energy, which is
\eqref{eq:hamiltonian} below; and \eqref{eq:euler-lagrange} becomes the pair
of first-order equations \eqref{eq:hamilton}. The gain is not the equations
themselves but the arena: the Hamiltonian formulation lives on phase space,
where the state is a \emph{point} whose evolution preserves volume, and that
is precisely the structure statistical mechanics needs to speak about
probability densities over states.

\section{Interacting systems: the Ising paradigm}\label{sec:ising}

Everything so far holds for arbitrary $H$; the physics (and the
difficulty) enters when the energy couples many degrees of freedom. The
canonical example is the \emph{Ising model} \citep{ising1925beitrag}:
binary spins $\sigma_i \in \{-1, +1\}$ on the nodes of a graph
$G = (V, \mathcal{E})$, with energy
\begin{equation}
  H(\sigma) \;=\; -\sum_{(i,j) \in \mathcal{E}} J_{ij}\, \sigma_i \sigma_j
  \;-\; \sum_{i \in V} h_i\, \sigma_i ,
  \label{eq:ising}
\end{equation}
where $J_{ij}$ are pairwise couplings and $h_i$ external fields. The
Gibbs distribution $p(\sigma) \propto e^{-\beta H(\sigma)}$ makes
neighbouring spins prefer alignment when $J_{ij} > 0$. Three lessons from
this model shape the rest of the thesis.

\paragraph{(i) The partition function is the hard part.}
$Z = \sum_\sigma e^{-\beta H(\sigma)}$ runs over $2^{|V|}$
configurations. On two-dimensional lattices Onsager computed it exactly
\citep{onsager1944crystal}; on general graphs the computation is
\#P-hard, and all of approximate inference (Monte Carlo sampling,
variational methods, message passing) exists because of this wall. On
\emph{trees}, however, $Z$ and all marginals are computable exactly in
time linear in the system size by the transfer-matrix recursion, the same
recursion that Chapter~\ref{ch:graphs} presents as belief propagation.
This tractable island is precisely where our research model lives.

\paragraph{(ii) Collective order emerges.} For $\beta$ beyond a critical
value the two-dimensional model magnetises spontaneously: a macroscopic
fraction of spins align even at $h = 0$, and the symmetric state splits
into two ergodic components. This spontaneous symmetry breaking is the
archetype of a \emph{phase transition}, and its dynamical counterpart, a
generation trajectory committing to one mode of a multimodal
distribution, is exactly the \emph{speciation} phenomenon of diffusion
models \citep{biroli2024dynamical, raya2023spontaneous}.

\paragraph{(iii) Disorder creates complexity.} Making the couplings
$J_{ij}$ random (of both signs) yields spin glasses, and, as the next
sections recount, spin glasses yielded neural networks. This is where the
physics of this chapter stops being background and becomes the
literature this thesis belongs to.

\section{Disorder, frustration, and spin glasses}\label{sec:spinglass}

In a ferromagnet all couplings agree, and the two ground states (all
spins up, all spins down) are found by inspection. Draw the couplings
$J_{ij}$ at random with both signs \citep{edwards1975theory,
sherrington1975solvable} and this transparency collapses. A triangle with
two positive bonds and one negative bond already cannot satisfy all
three: whatever the spins do, at least one bond pays. This is
\emph{frustration}, and its large-$N$ consequence is a free-energy
landscape shattered into exponentially many valleys separated by
extensive barriers: the system has many near-equally-good ways of being,
none of them readable off the couplings. The mean-field theory of this
phenomenon, developed through the replica and cavity methods
\citep{mezard1987spin}, and the Thouless--Anderson--Palmer (TAP)
self-consistency equations \citep{thouless1977solution}, took a decade to
control and earned Parisi the 2021 Nobel Prize in Physics.

Two exports from this episode matter here. First, a language: rugged
landscapes, metastable states, order parameters that measure similarity
rather than magnetisation. Machine learning speaks this language every
time it discusses loss surfaces and local minima. Second, a technology:
the cavity and TAP equations are recursions on \emph{local} quantities
that solve a \emph{global} inference problem, and they are the direct
ancestors of the message-passing algorithms \citep{mezard2009information,
zdeborova2016statistical, krzakala2024statphys} that
Chapter~\ref{ch:graphs} develops and that the research chapters deploy on
diffusion scores. The connection between disordered magnets and these
algorithms is historical descent, not analogy.

\section{The Hopfield network: memory as a phase of matter}
\label{sec:hopfield}

\citet{hopfield1982neural} asked a physicist's question about memory: can
a system of simple binary units \emph{store} and \emph{retrieve} patterns
as a collective effect, the way a magnet stores a magnetisation? His
construction is an Ising system \eqref{eq:ising} in disguise. Take $N$
binary neurons $\sigma_i \in \{-1,+1\}$ and $P$ patterns
$\xi^\mu \in \{-1,+1\}^N$ to be memorised, and wire the couplings by the
\emph{Hebb rule}, ``neurons that fire together wire together''
\citep{hebb1949organization}:
\begin{equation}
  J_{ij} \;=\; \frac{1}{N} \sum_{\mu=1}^{P} \xi_i^\mu \xi_j^\mu ,
  \qquad
  H(\sigma) \;=\; -\frac{1}{2}\sum_{i \neq j} J_{ij}\, \sigma_i \sigma_j .
  \label{eq:hopfield}
\end{equation}
The dynamics is the simplest imaginable: pick a neuron, align it with its
local field, repeat,
\begin{equation}
  \sigma_i \;\leftarrow\; \operatorname{sign}\Big( \textstyle\sum_{j \neq i}
  J_{ij}\, \sigma_j \Big).
  \label{eq:hopfield-dynamics}
\end{equation}

\begin{toolbox}{The Hopfield dynamics descends the energy}
Let $h_i = \sum_{j \neq i} J_{ij} \sigma_j$ be the local field on neuron
$i$ and suppose the update \eqref{eq:hopfield-dynamics} flips it,
$\sigma_i' = -\sigma_i$. Only the terms of $H$ containing $\sigma_i$
change, so
\begin{equation*}
  \Delta H
  \;=\; H(\sigma') - H(\sigma)
  \;=\; -\,(\sigma_i' - \sigma_i)\, h_i
  \;=\; -\,2\,\sigma_i'\, h_i .
\end{equation*}
The update sets $\sigma_i' = \operatorname{sign}(h_i)$, hence
$\sigma_i' h_i = |h_i| \geq 0$ and $\Delta H \leq 0$, with strict
decrease whenever a flip actually occurs. Since $H$ is bounded below on a
finite configuration space, the dynamics must terminate in a local
minimum: \emph{every} initial condition falls into an attractor. Memory
retrieval is this fall: initialise at a corrupted pattern, descend, and
read off the fixed point. \hfill$\square$
\end{toolbox}

The patterns are, by construction, near the bottoms of the energy
\eqref{eq:hopfield}: retrieval works because each memory is an attractor
with a basin. How many memories fit is a genuine phase-transition
question, answered with spin-glass tools by \citet{amit1985storing}: for
random patterns with $P = \alpha N$, retrieval survives only below a
critical load $\alpha_c \approx 0.138$, beyond which the memories are
crushed into a useless glassy phase. In this model, memory is a phase
of matter in the technical sense: it is entered and left through sharp
transitions. The
single-pattern case $P=1$, known as the Mattis model
\citep{mattis1976solvable}, is gauge-equivalent to the uniform
ferromagnet: flipping $\sigma_i \to \xi_i \sigma_i$ maps the stored
pattern onto the all-up state. This observation is worth keeping in mind:
the Mattis model returns in Section~\ref{sec:cascade} as the exactly
solvable teacher in a modern analysis of how energy-based models learn.

In 2024 the Nobel Prize in Physics went to John Hopfield and Geoffrey
Hinton ``for foundational discoveries and inventions that enable machine
learning with artificial neural networks'' \citep{nobel2024physics}:
formal recognition that the lineage traced in this section is physics,
not just inspired by it.

\section{From Hopfield to Boltzmann machines: learning enters}
\label{sec:boltzmann-machines}

The Hopfield network stores what its designer wires into it. The decisive
step towards machine \emph{learning} was to make the couplings themselves
the unknowns. Promote the zero-temperature dynamics
\eqref{eq:hopfield-dynamics} to stochastic updates at temperature
$T = 1/\beta$ (each neuron aligns with its field only with probability
given by the Boltzmann factor); the stationary law of this dynamics is
exactly the Gibbs measure $p(\sigma) \propto e^{-\beta H(\sigma)}$ of the
Ising energy. A \emph{Boltzmann machine} \citep{ackley1985learning} is
this object run backwards: fix a dataset, and ask for the couplings $J$
that make the Gibbs measure reproduce it. Gradient ascent on the
log-likelihood gives a rule of striking symmetry,
\begin{equation}
  \frac{\partial}{\partial J_{ij}} \,\E_{\text{data}}\big[\log p(\sigma)\big]
  \;=\; \beta\,\Big(
  \underbrace{\langle \sigma_i \sigma_j \rangle_{\text{data}}}_{\text{positive phase}}
  \;-\;
  \underbrace{\langle \sigma_i \sigma_j \rangle_{\text{model}}}_{\text{negative phase}}
  \Big),
  \label{eq:bm-learning}
\end{equation}
and learning stops when the model's correlations match the data's. This
is the maximum-entropy principle of Toolbox~\ref{tb:maxent} turned into
an algorithm: matching prescribed moments with an exponential-family
measure. The rule also contains, in miniature, the central obstruction of
the field. The positive phase is an average over data, which is cheap.
The negative phase is an average over the \emph{model's} Gibbs measure: a
partition-function computation, the wall of Section~\ref{sec:ising},
approachable only through sampling
(Section~\ref{sec:mcmc}).

The \emph{restricted} Boltzmann machine (RBM) tames this cost by
architecture \citep{smolensky1986information}: split the units into a
visible layer $v$ and a hidden layer $h$ with couplings only across
layers, $E(v,h) = -v^\T W h - b^\T v - c^\T h$. On this bipartite graph
each layer is conditionally independent given the other, so Gibbs
sampling alternates two exact, fully parallel steps, and the
contrastive-divergence approximation to \eqref{eq:bm-learning} made
training practical \citep{hinton2002training}. Stacked RBMs, the deep
belief networks, delivered the first working deep architectures and
relaunched neural networks as a field \citep{hinton2006fast}. The lineage
then split. The supervised branch, powered by backpropagation
\citep{rumelhart1986learning}, scaled through convolutional networks to
the modern deep-learning era \citep{krizhevsky2012imagenet,
lecun2015deep} and largely left equilibrium physics behind. The
\emph{generative} branch, the one this thesis lives in, kept the Gibbsian
core: a model is an energy, learning is moment matching, sampling is
dynamics. Chapter~\ref{ch:diffusion} follows that branch from
energy-based models to diffusion.

\section{The statistical mechanics of learning}\label{sec:statmech-learning}

Physics did not stop contributing machines; it also built the theory of
their behaviour. \citet{gardner1988space} showed that learning itself can
be treated as an equilibrium problem: fix a set of examples and compute,
by replica methods, the volume of coupling space compatible with them;
capacity, generalisation, and their transitions then come out as
thermodynamics in \emph{weight} space. This programme, extended through
teacher--student settings, matured into the modern statistical physics of
inference \citep{zdeborova2016statistical}, whose central lesson recurs
throughout this thesis: high-dimensional inference problems have phases,
and algorithms fail or succeed at sharp boundaries.

Two strands of this literature frame our research questions directly.
The first is algorithmic: \citet{mezard2017meanfield} derived the
mean-field message-passing (TAP) equations of the Hopfield model and its
generalisations, exhibiting the equivalence between Hopfield networks and
RBMs and showing that retrieval can be organised as a message-passing
computation. It is the same machinery, transplanted from memories to
diffusion scores, that the research chapters run on our model. The second
is dynamical: training itself undergoes phase transitions. In linear
networks, gradient descent learns the singular modes of the data
covariance sequentially \citep{saxe2014exact}; in RBMs, the weight
matrix's singular vectors align with the data's principal components in
the early phase of learning \citep{decelle2017spectral}, and hidden units
develop compositional receptive fields through condensation transitions
\citep{tubiana2017emergence}. The sharpest result in this line,
\citet{bachtis2024cascade}, examined in depth in
Section~\ref{sec:cascade}, shows that RBM training is an entire
\emph{cascade} of second-order phase transitions, one per learned
feature. Learning, like memory, turns out to be a sequence of symmetry
breakings.

The order of a transition is the order of the derivative of the free energy
that first misbehaves, and the distinction matters for reading these results.
At a \emph{first-order} transition the first derivative jumps: the order
parameter is discontinuous, latent heat is absorbed, the two phases coexist at
the transition, and correlation lengths stay finite. At a \emph{second-order}
transition the first derivative is continuous and the second diverges: the
order parameter grows continuously from zero, the susceptibility
\eqref{eq:static-response} and the correlation length diverge, and observables
acquire power laws. The cascade above is of the second kind, which is why it is
detected through susceptibilities rather than through jumps. Both notions
require the thermodynamic limit; in the finite models of
Chapters~\ref{ch:gaussian} to~\ref{ch:ring} a correlation length that grows
until it reaches the size of the system and then stops is a finite-size
ceiling, not a divergence.

\section{What this thesis takes from physics}

A \emph{complex system} is, operationally, a system whose macroscopic
behaviour is not readable off its microscopic rules: interactions
generate collective phenomena (order, phases, transitions, slow dynamics)
that must be derived, not postulated. Statistical mechanics is the
toolbox for that derivation, and after this chapter's story its objects
map one-to-one onto the objects of modern generative modelling:
\begin{center}
\begin{tabular}{ll}
\toprule
\textbf{Statistical mechanics} & \textbf{Generative modelling} \\
\midrule
microstate $x$ & data point / latent configuration \\
energy $H(x)$ & negative log-density $-\log p(x)$ (up to $Z$) \\
Boltzmann distribution $e^{-\beta H}/Z$ & model distribution \\
partition function $Z$ & normalising constant / evidence \\
free energy $F$ & negative ELBO (up to constants) \\
Hebbian couplings, attractors & stored memories, retrieval \\
moment matching \eqref{eq:bm-learning} & maximum-likelihood training \\
Langevin dynamics & sampling / reverse diffusion \\
phase transition & feature condensation, speciation \\
cavity / TAP equations & belief propagation / AMP \\
\bottomrule
\end{tabular}
\end{center}
The next chapter supplies the probabilistic dynamics, Markov chains and
their continuous-time limits, that turn the static ensembles of this
chapter into processes that move probability around, which is what a
diffusion model ultimately is.
